%=====================================================================
\section{The environment: \code{rl\_f4\_env.py}}
\label{sec:env}

The environment is the ``game'' the policy plays. It answers three questions at every step: \emph{which decisions are allowed?}, \emph{what does each cost?}, and \emph{what does the partial plan look like after the decision?} Everything is written with PyTorch tensors so that many episodes (many instances, or many tries on one instance) run side by side; the first dimension of every tensor is the \emph{row} $b$ of the batch.

\subsection{\code{Batch}: the static data as tensors}

\src{../code/rl_f4_env.py}{53}{65}

\begin{itemize}
  \item A batch holds $B$ instances with the same $n$, stations, guns and vans (line 59 checks it), so that they fit in one tensor.
  \item Lines 62--63 number the guns: station 1 has guns 0 and 1, station 2 has guns 2 and 3, and so on. \code{gun\_station[g]} is $f(g)$; with 12 stations of 2 guns, $G=24$.
  \item Line 65: the resources are the $G$ guns followed by the $M$ vans, $K=G+M=29$. A resource index $k<G$ is a gun, $k\ge G$ is van $k-G$.
\end{itemize}

Lines 71--87 copy every number of \code{J} into tensors of shape $(B,n)$, $(B,n,F)$, $(B,n{+}1,n{+}1)$ or $(B,)$, in 64-bit floating point (so feasibility decisions agree exactly with the checker). Lines 88--99 then compute quantities that never change during an episode:

\src{../code/rl_f4_env.py}{88}{100}

\begin{itemize}
  \item \code{reach\_g}$[i,g]=\big[d_{i f(g)}\le R_i\big]$: car $i$ can reach gun $g$'s station;
  \item \code{arrive\_g}$[i,g]=r_i+t_{i f(g)}$: when car $i$ would arrive there;
  \item cost of gun $g$ for car $i$ without any queue:
  \[
  \text{\code{gun\_cost0}}[i,g]=\mu_i\big(c^{EV}d_{if(g)}+p^\Fcs E_i\big)+\pi^\Fcs E_i+c^D_it_{if(g)}+c^C_i\tau^\Fcs_i ;
  \]
  \item cost of a van for car $i$ without wait and without the van's driving:
  \[
  \text{\code{van\_cost0}}[i]=\mu_ip^\Mcs E_i+\pi^\Mcs\frac{E_i}{\eta}+c^V_i\tau^\Mcs_i ;
  \]
  \item \code{d\_back}$[i]=d_{i0}$, \code{t\_back}$[i]=t_{i0}$: from car $i$ back to the depot.
\end{itemize}
\code{gun\_cost0} and \code{van\_cost0} are the parts of the cost increase that depend only on the car and the resource type; the parts that depend on the current state (queue, van wait, van position, whether the van is new) are added in \code{options()}.

\subsection{\code{Env}: the state}

\src{../code/rl_f4_env.py}{126}{157}

{\small
\begin{longtable}{@{}>{\raggedright\arraybackslash}p{2.6cm}>{\raggedright\arraybackslash}p{1.6cm}>{\raggedright\arraybackslash}p{2.4cm}>{\raggedright\arraybackslash}p{8.6cm}@{}}
\toprule
tensor & shape & start value & meaning \\
\midrule
\endhead
\code{done} & $(B,n)$ & false & car $i$ is placed or dropped \\
\code{dropped} & $(B,n)$ & false & car $i$ could not be placed (no allowed resource left) \\
\code{gun\_free} & $(B,G)$ & 0 & $a_g$: the minute gun $g$ becomes free (completion of its last car) \\
\code{van\_pos} & $(B,M)$ & 0 & $\ell_m$: node where van $m$ is (0 = depot, $i+1$ = car $i$) \\
\code{van\_free} & $(B,M)$ & 0 & $a_m$: the minute van $m$ finishes its last car \\
\code{van\_batt} & $(B,M)$ & $B$ & $b_m$: battery left (kWh) \\
\code{van\_cnt} & $(B,M)$ & 0 & $\kappa_m$: cars served so far \\
\code{cost} & $(B,)$ & 0 & sum of the cost increases so far ($=Z$ of the partial plan) \\
\code{res}, \code{S}, \code{C}, \code{order} & $(B,n)$ & $-1$, 0, 0, $-1$ & the plan: resource, start, completion and step number of each car \\
\code{steps} & int & 0 & decisions taken \\
\bottomrule
\end{longtable}}

This is the full \textbf{state} of the decision process:
\[
s_t=\Big(P_t,\ D_t,\ (a_g)_{g\in G},\ (\ell_m,a_m,b_m,\kappa_m)_{m\in M}\Big),
\]
with $P_t$ the placed and $D_t$ the dropped cars. Everything the future depends on is in it: the cost of every remaining decision can be computed from $s_t$ alone (this is the \emph{Markov property}).

Two extra quantities are prepared once:
\begin{itemize}
  \item \code{same\_earlier[g,g']} (line 149) is true when $g'<g$ and both guns belong to the same station. It is used for the symmetry rule below.
  \item \code{cost\_scale} (lines 152--157), written $\sigma$: the average, over the cars, of their cheapest allowed option at the start. It is a ``typical cost of one car'' for this instance, used to scale features and the penalty for unplaced cars:
  \[
  \sigma=\max\Big(1,\ \frac{1}{|\{i:\exists k\}|}\sum_{i}\min_{k\text{ allowed}}\Delta Z(i,k)\Big).
  \]
\end{itemize}

\subsection{The action}

An \textbf{action} is a pair $(i,k)$: car $i$ (not yet done) goes to resource $k$. It is appended to the \emph{end} of $k$'s sequence: the end of gun $k$'s queue, or the next stop of van $k-G$'s route. The policy chooses the pair in two steps, first $i$ and then $k$ (Section~\ref{sec:network}); the environment does not care how it was chosen.

\subsection{\code{options()}: what every possible action would do}
\label{sec:options}

\code{options()} computes, for \emph{every} unplaced car $i$ and \emph{every} resource $k$ at once, the start, wait, finish, whether the pair is allowed (the \emph{mask}), and the cost increase $\Delta Z(i,k)$. The result has shape $(B,n,K)$.

\src{../code/rl_f4_env.py}{160}{195}

\paragraph{Guns (lines 165--172).} If car $i$ is sent to gun $g$:
\begin{align*}
S_i&=\max\big(r_i+t_{if(g)},\ a_g\big) &&\text{arrive, or wait until the gun is free (line 166)}\\
q^\Fcs_i&=S_i-r_i-t_{if(g)} &&\text{queue (line 167)}\\
C_i&=S_i+\tau^\Fcs_i &&\text{finish (line 168)}
\end{align*}
The pair is \textbf{allowed} (line 170) if all four hold:
\[
d_{if(g)}\le R_i,\qquad q^\Fcs_i\le W_i,\qquad C_i\le T,
\]
and no earlier gun of the same station is free at the same minute.
The last condition (lines 169, 171) is the \emph{symmetry rule}: two idle guns of one station are identical, so only the first is offered. The cost increase (line 172) is
\[
\Delta Z(i,g)=\text{\code{gun\_cost0}}[i,g]+c^Q_i\,q^\Fcs_i .
\]

\paragraph{Vans (lines 174--191).} If car $i$ is the next stop of van $m$, which is at node $\ell_m$ and free at $a_m$:
\begin{align*}
S_i&=\max\big(r_i,\ a_m+t_{\ell_m i}\big) &&\text{the van drives over; the car must have requested (line 179)}\\
q^\Mcs_i&=S_i-r_i &&\text{wait for the van (line 180)}\\
C_i&=S_i+\tau^\Mcs_i &&\text{finish (line 181)}\\
\text{need}_{im}&=\frac{E_i}{\eta}+\beta\big(d_{\ell_m i}+d_{i0}\big) &&\text{energy to charge $i$, drive there and get home (line 182)}
\end{align*}
The pair is \textbf{allowed} (lines 185--187) if all four hold:
\[
\kappa_m<Q,\qquad \text{need}_{im}\le b_m,\qquad C_i+t_{i0}\le T,
\]
and the van is already in use ($\kappa_m>0$) or is the first unused van.
The last condition is again a symmetry rule: unused vans are identical (all at the depot, full battery), so only the first is offered (lines 183--184). The cost increase (lines 188--191) is
\[
\Delta Z(i,m)=\text{\code{van\_cost0}}[i]+c^\Mcs\big(d_{\ell_m i}+d_{i0}-d_{\ell_m0}\big)+c^\Fcs\,[\kappa_m=0]+c^V_i\,q^\Mcs_i .
\]
The van-km term is the \emph{extra} driving: before, the van's route ended $\ell_m\to0$; now it ends $\ell_m\to i\to0$. The dispatch cost $c^\Fcs$ is charged when a van gets its first car.

\paragraph{Line 192.} Cars that are already done get no allowed resource. The returned dictionary keeps the full $(B,n,K)$ tensors (\code{feas}, \code{finish}, \code{dz}) and the gun and van parts of start and wait separately (they are read later only for the chosen car, which saves memory).

\subsection{\code{drop\_dead()} and \code{step()}: what happens after a decision}

\src{../code/rl_f4_env.py}{201}{235}

\textbf{\code{drop\_dead}} (lines 201--207): a car with no allowed resource is marked done \emph{and} dropped. It can never become placeable again (Proposition~\ref{prop:dead}), so waiting would only waste steps.

\textbf{\code{step}} (lines 209--235) applies action $(i,k)$ in every active row:
\begin{itemize}
  \item lines 216--218: read the start and finish computed by \code{options()};
  \item line 219: \code{cost} $\leftarrow$ \code{cost} $+\Delta Z(i,k)$;
  \item lines 220--223: record the plan (resource, $S_i$, $C_i$, step number) and mark the car done;
  \item gun (line 226): $a_g\leftarrow C_i$;
  \item van (lines 229--234):
  \[
  b_m\leftarrow b_m-\frac{E_i}{\eta}-\beta\,d_{\ell_m i},\qquad \ell_m\leftarrow i,\qquad a_m\leftarrow C_i,\qquad \kappa_m\leftarrow\kappa_m+1 .
  \]
  The way home is \emph{not} subtracted from the battery yet; it was only checked in the mask, and it will be checked again from the new position for the next car.
\end{itemize}
The environment is \textbf{deterministic}: the same actions always give the same plan.

\textbf{End of an episode} (\code{finished()}, line 237): when every car is done (placed or dropped).

\subsection{What the network sees: the features}
\label{sec:features}

The network never sees the raw tensors; \code{rl\_f4\_env.py} turns the state into scaled numbers.

\paragraph{Static features} (lines 245--274), computed once per instance:

\src{../code/rl_f4_env.py}{263}{273}

{\small
\begin{longtable}{@{}r>{\raggedright\arraybackslash}p{6.6cm}>{\raggedright\arraybackslash}p{8.0cm}@{}}
\toprule
\# & car feature & formula \\
\midrule
\endhead
1--2 & position relative to the depot & $(x_i-x_0)/\text{scale}$, one of 8 map symmetries \\
3 & request time & $r_i/T$ \\
4--5 & charging times & $\tau^\Fcs_i/T$, $\tau^\Mcs_i/T$ \\
6 & queue tolerance & $W_i/T$ \\
7 & income & $\theta_i$ \\
8 & energy & $E_i/B$ \\
9 & stranded & $1$ if no station is within range \\
10 & share of reachable stations & $|F_i|/F$ \\
11 & time to the nearest reachable station & $\min_{f\in F_i}t_{if}/T$ \\
12 & time back to the depot & $t_{i0}/T$ \\
13 & range relative to the city & $R_i/\max d_{if}$ \\
14 & cheapest gun cost & $\min_g\text{\code{gun\_cost0}}[i,g]/\sigma$ \\
15 & van cost & $\text{\code{van\_cost0}}[i]/\sigma$ \\
16 & how late a van may come & $(T-r_i-\tau^\Mcs_i-t_{i0})/T$ \\
17 & \textbf{commercial} (new in f4) & $s_i\in\{0,1\}$ \\
18 & \textbf{money sensitivity} (new) & $\mu_i/2$ \\
19 & \textbf{cost of queueing one charge} (new) & $c^Q_i\tau^\Fcs_i/\sigma$ \\
\bottomrule
\end{longtable}}
Each station has 4 features (position, guns/2, share of cars that can reach it); the depot has 4 (seats per car $QM/n$, guns per car $G/n$, $\log T/8$, $n/100$).

\paragraph{Dynamic features}, recomputed at every step from \code{options()} and the state:
\begin{itemize}
  \item \textbf{global} (4, lines 276--281): share of cars done, mean gun busy time $/T$, share of vans in use, mean battery left $/B$;
  \item \textbf{per car} $\varphi_i$ (7, lines 283--314): share of resources still allowed; cheapest $\Delta Z/\sigma$; \emph{regret} = (second-cheapest $-$ cheapest)$/\sigma$; earliest finish $/T$; (cheapest gun $-$ cheapest van)$/\sigma$; request time relative to the earliest open request; place in the queue of open requests (0 = next, 1 = last; ties broken by higher $c^Q$);
  \item \textbf{per resource} $\rho_k$ (6, lines 316--325): free time $/T$, van flag, battery left $/B$, seats used $/Q$, in use, time back to the depot $/T$ (van-only entries are 0 for guns);
  \item \textbf{per pair} (6, lines 327--344), only for the chosen car $i$: $\Delta Z(i,k)/\sigma$, its excess over the cheapest option, wait $/T$, finish $/T$, battery left after serving $i$, van flag; all zero for masked resources.
\end{itemize}

\subsection{\code{GreedyRule}: the simple rule}

\src{../code/rl_f4_env.py}{348}{365}

The order is built with three stable sorts (lines 356--358): first by $c^Q$ (highest first), then by request time, then stranded cars first, so the final order is \emph{stranded first, then earliest request, then the customer who hates queueing most}. At each step the first car in that order that is not done takes its cheapest allowed resource (lines 363--365). This is the MILP's own starting plan (\code{greedy\_f4}); in the RL it is the teacher of the warm-up and the yardstick of the validation.

\subsection{\code{plan\_of}: from an episode to a plan}

\src{../code/rl_f4_env.py}{423}{439}

Lines 423--438 turn the recorded resources into the plan format shared with the MILP: ``FCS$f$-gun$g$'' for station users, ``van$v$'' for van users (vans in use numbered $1,2,\dots$ like the MILP), and each van's route $[0,\dots,0]$ in the order the cars were added. Lines 395--422 (not shown) handle dropped cars: they are appended to whichever resource finishes them earliest, \emph{ignoring} the rule they break, so the plan is complete; the checker then lists the violation and the plan is reported as \code{invalid\_plan}.

\inshort{The environment keeps the partial plan (when each gun is free; where each van is, when it is free, how much battery and how many seats it has left). For every possible next decision it computes the start time, whether any rule is broken, and the exact cost increase. Applying a decision updates the plan and adds the cost.}
